\documentclass[12pt,a4paper]{article}

\usepackage[T2A]{fontenc}
\usepackage[utf8]{inputenc}
\usepackage[russian]{babel}
\usepackage{amsmath,amssymb,amsthm,mathtools}
\usepackage{enumitem}
\usepackage{listings}
\usepackage{xcolor}
\usepackage{geometry}

\geometry{margin=2.2cm}

\newtheorem{corollary}{Следствие}
\newtheorem{definition}{Определение}
\newtheorem{theorem}{Теорема}
\newtheorem{lemma}{Лемма}
\newtheorem{example}{Пример}
\newtheorem{remark}{Замечание}

\lstdefinestyle{fpstyle}{
  basicstyle=\ttfamily\small,
  keywordstyle=\color{blue},
  commentstyle=\color{gray},
  stringstyle=\color{teal},
  showstringspaces=false,
  frame=single,
  breaklines=true,
  tabsize=2,
  columns=fullflexible
  literate=
    {А}{{\CYRA}}1
    {Б}{{\CYRB}}1
    {В}{{\CYRV}}1
    {Г}{{\CYRG}}1
    {Д}{{\CYRD}}1
    {Е}{{\CYRE}}1
    {Ё}{{\CYRYO}}1
    {Ж}{{\CYRZH}}1
    {З}{{\CYRZ}}1
    {И}{{\CYRI}}1
    {Й}{{\CYRISHRT}}1
    {К}{{\CYRK}}1
    {Л}{{\CYRL}}1
    {М}{{\CYRM}}1
    {Н}{{\CYRN}}1
    {О}{{\CYRO}}1
    {П}{{\CYRP}}1
    {Р}{{\CYRR}}1
    {С}{{\CYRS}}1
    {Т}{{\CYRT}}1
    {У}{{\CYRU}}1
    {Ф}{{\CYRF}}1
    {Х}{{\CYRH}}1
    {Ц}{{\CYRC}}1
    {Ч}{{\CYRCH}}1
    {Ш}{{\CYRSH}}1
    {Щ}{{\CYRSHCH}}1
    {Ъ}{{\CYRHRDSN}}1
    {Ы}{{\CYRERY}}1
    {Ь}{{\CYRSFTSN}}1
    {Э}{{\CYREREV}}1
    {Ю}{{\CYRYU}}1
    {Я}{{\CYRYA}}1
    {а}{{\cyra}}1
    {б}{{\cyrb}}1
    {в}{{\cyrv}}1
    {г}{{\cyrg}}1
    {д}{{\cyrd}}1
    {е}{{\cyre}}1
    {ё}{{\cyryo}}1
    {ж}{{\cyrzh}}1
    {з}{{\cyrz}}1
    {и}{{\cyri}}1
    {й}{{\cyrishrt}}1
    {к}{{\cyrk}}1
    {л}{{\cyrl}}1
    {м}{{\cyrm}}1
    {н}{{\cyrn}}1
    {о}{{\cyro}}1
    {п}{{\cyrp}}1
    {р}{{\cyrr}}1
    {с}{{\cyrs}}1
    {т}{{\cyrt}}1
    {у}{{\cyru}}1
    {ф}{{\cyrf}}1
    {х}{{\cyrh}}1
    {ц}{{\cyrc}}1
    {ч}{{\cyrch}}1
    {ш}{{\cyrsh}}1
    {щ}{{\cyrshch}}1
    {ъ}{{\cyrhrdsn}}1
    {ы}{{\cyrery}}1
    {ь}{{\cyrsftsn}}1
    {э}{{\cyrerev}}1
    {ю}{{\cyryu}}1
    {я}{{\cyrya}}1
}

\lstset{style=fpstyle}

\title{Билет №48 \\ \small Достоинства и недостатки автоматической генерации кода системы по модели. (СпБПУ)}
\author{Сериков Владислав Александрович}
\date{13 мая 2026}

\begin{document}

\maketitle
\tableofcontents
\newpage

\section{Автоматическая генерация кода системы по модели}

\subsection{Основные понятия}

\textbf{Модель} --- формализованное представление системы, отражающее существенные свойства предметной области, архитектуры, поведения или реализации и абстрагирующееся от несущественных деталей.

\textbf{Автоматическая генерация кода по модели} --- процесс получения исходного кода, конфигурационных файлов, схем баз данных, интерфейсов, тестов или иной программной артефактной базы на основании формальной или полуформальной модели системы.

Идея модельно-ориентированной разработки состоит в том, что основным артефактом разработки становится не только исходный код, но и модель, из которой код может быть частично или полностью получен автоматически.

\subsection{Место генерации кода в модельно-ориентированной разработке}

В модельно-ориентированной разработке обычно выделяют следующие уровни:

\begin{enumerate}
    \item \textbf{CIM} --- computation independent model, модель требований и предметной области, независимая от вычислительной реализации.
    \item \textbf{PIM} --- platform independent model, платформенно-независимая модель логики системы.
    \item \textbf{PSM} --- platform specific model, платформенно-зависимая модель, учитывающая конкретную технологию, язык программирования, фреймворк, СУБД.
    \item \textbf{Code} --- исходный код и сопутствующие артефакты.
\end{enumerate}

Типичная цепочка преобразований имеет вид:
\[
    CIM \longrightarrow PIM \longrightarrow PSM \longrightarrow Code.
\]

На практике возможны и более простые схемы:
\[
    UML\text{-диаграмма классов} \longrightarrow Java\text{-классы},
\]
\[
    ER\text{-модель} \longrightarrow SQL\text{-схема},
\]
\[
    BPMN\text{-модель процесса} \longrightarrow исполняемый workflow.
\]

\subsection{Виды генерации кода}

\subsubsection{Полная генерация}

При полной генерации весь исполняемый код получается из модели автоматически.

\textbf{Пример:}
\[
    \text{модель конечного автомата} \longrightarrow \text{код обработчика состояний}.
\]

Такой подход хорошо работает в ограниченных предметных областях: встроенные системы, телекоммуникационные протоколы, генерация CRUD-приложений, генерация конфигураций.

\subsubsection{Частичная генерация}

При частичной генерации автоматически создаётся только часть системы:

\begin{itemize}
    \item классы предметной области;
    \item интерфейсы;
    \item DTO-объекты;
    \item API-контроллеры;
    \item схемы БД;
    \item тестовые шаблоны;
    \item конфигурации;
    \item документация.
\end{itemize}

Оставшаяся часть дорабатывается вручную.

\subsubsection{Генерация каркаса}

Генератор создаёт шаблонную структуру проекта:

\begin{itemize}
    \item директории;
    \item базовые классы;
    \item заглушки методов;
    \item интерфейсы;
    \item конфигурационные файлы.
\end{itemize}

Разработчик затем вручную реализует бизнес-логику.

\subsubsection{Генерация с сохранением пользовательского кода}

В этом случае генератор может запускаться многократно, не уничтожая ручные изменения.

Обычно используются следующие техники:

\begin{enumerate}
    \item protected regions --- специальные области в сгенерированном коде;
    \item наследование --- генератор создаёт базовый класс, пользователь расширяет его;
    \item partial classes --- разделение класса на сгенерированную и ручную части;
    \item композиция --- сгенерированный код вызывает пользовательские расширения;
    \item шаблон ``generation gap''.
\end{enumerate}

\subsection{Общая схема автоматической генерации кода}

Пусть имеется модель:
\[
    M = (E, R, A),
\]
где:
\begin{itemize}
    \item $E$ --- множество элементов модели;
    \item $R$ --- множество отношений между элементами;
    \item $A$ --- множество атрибутов элементов.
\end{itemize}

Генератор кода можно рассматривать как отображение:
\[
    G: M \rightarrow C,
\]
где $C$ --- множество программных артефактов.

Например:
\[
    G(\text{ClassDiagram}) = \{\text{Java-классы}, \text{интерфейсы}, \text{DAO}, \text{тесты}\}.
\]

\subsection{Алгоритм генерации кода по модели}

\subsubsection{Общий алгоритм}

\begin{enumerate}
    \item \textbf{Построить или загрузить модель.}

    Модель может быть создана в UML-редакторе, DSL-редакторе, CASE-средстве или описана текстово.

    \item \textbf{Проверить корректность модели.}

    Проверяются:
    \begin{itemize}
        \item синтаксическая корректность;
        \item соответствие метамодели;
        \item отсутствие противоречий;
        \item выполнение ограничений.
    \end{itemize}

    \item \textbf{Выполнить преобразование модели.}

    При необходимости модель преобразуется из платформенно-независимой в платформенно-зависимую:
    \[
        PIM \rightarrow PSM.
    \]

    \item \textbf{Выбрать шаблоны генерации.}

    Например:
    \begin{itemize}
        \item шаблон класса;
        \item шаблон интерфейса;
        \item шаблон репозитория;
        \item шаблон REST-контроллера;
        \item шаблон SQL-таблицы.
    \end{itemize}

    \item \textbf{Выполнить подстановку элементов модели в шаблоны.}

    Для каждого элемента модели создаются соответствующие программные конструкции.

    \item \textbf{Сгенерировать файлы.}

    На этом шаге создаются исходные файлы, конфигурации, схемы, тесты.

    \item \textbf{Отформатировать и проверить сгенерированный код.}

    Выполняются:
    \begin{itemize}
        \item форматирование;
        \item статический анализ;
        \item компиляция;
        \item запуск тестов.
    \end{itemize}

    \item \textbf{Интегрировать сгенерированный код в проект.}

    Код добавляется в репозиторий или артефактную сборку.
\end{enumerate}

\subsubsection{Минимальный пример}

Пусть имеется простая модель класса:

\[
\begin{array}{|c|}
\hline
\text{User} \\
\hline
id: Integer \\
name: String \\
email: String \\
\hline
\end{array}
\]

Эта модель может быть преобразована в Java-класс:

\begin{verbatim}
public class User {
    private Integer id;
    private String name;
    private String email;

    public Integer getId() {
        return id;
    }

    public void setId(Integer id) {
        this.id = id;
    }

    public String getName() {
        return name;
    }

    public void setName(String name) {
        this.name = name;
    }

    public String getEmail() {
        return email;
    }

    public void setEmail(String email) {
        this.email = email;
    }
}
\end{verbatim}

Правило генерации можно записать следующим образом:

\[
    \text{Class}(name, attributes) \mapsto \text{JavaClass}(name, fields, getters, setters).
\]

Если атрибут модели имеет вид:
\[
    a = (name, type),
\]
то для него генерируется поле:
\[
    private\ type\ name;
\]
и пара методов доступа:
\[
    getName(), \quad setName(type\ value).
\]

\subsection{Шаблонная генерация}

Наиболее распространённый способ генерации кода основан на шаблонах.

\textbf{Шаблон} --- текстовая заготовка, содержащая фиксированные фрагменты кода и переменные части, заполняемые данными из модели.

Пример шаблона:

\begin{verbatim}
public class ${className} {
<#list attributes as attr>
    private ${attr.type} ${attr.name};
</#list>
}
\end{verbatim}

Если в модель передать:

\begin{verbatim}
className = "Product"
attributes = [
    {type = "Long", name = "id"},
    {type = "String", name = "title"},
    {type = "Double", name = "price"}
]
\end{verbatim}

то результатом будет:

\begin{verbatim}
public class Product {
    private Long id;
    private String title;
    private Double price;
}
\end{verbatim}

\subsection{Достоинства автоматической генерации кода}

\subsubsection{Повышение производительности разработки}

Одно из главных достоинств генерации кода состоит в сокращении объёма ручного программирования.

Если в системе много однотипных элементов, например:

\begin{itemize}
    \item сущностей;
    \item DTO;
    \item CRUD-контроллеров;
    \item репозиториев;
    \item SQL-таблиц;
    \item форм пользовательского интерфейса;
\end{itemize}

то генератор позволяет создавать их автоматически.

Например, для $n$ сущностей ручная реализация CRUD-слоя может требовать порядка:
\[
    O(n \cdot k)
\]
ручных действий, где $k$ --- среднее число шаблонных операций на одну сущность.

При генерации человек описывает только модель, а повторяющиеся операции выполняются автоматически.

\subsubsection{Снижение количества рутинных ошибок}

Ручное написание однотипного кода часто приводит к ошибкам:

\begin{itemize}
    \item опечаткам в именах;
    \item несогласованности типов;
    \item пропущенным методам;
    \item ошибкам копирования;
    \item несоответствию между моделью данных и кодом.
\end{itemize}

Автоматическая генерация снижает вероятность таких ошибок, поскольку одинаковые правила применяются систематически.

Если правило генерации корректно, то все однотипные фрагменты кода будут сформированы единообразно.

\subsubsection{Единообразие архитектуры}

Генератор навязывает проекту общую структуру:

\begin{itemize}
    \item единый стиль именования;
    \item одинаковое расположение классов;
    \item единый способ обработки ошибок;
    \item одинаковые шаблоны контроллеров;
    \item согласованную структуру модулей.
\end{itemize}

Это особенно важно в больших командах, где разные разработчики могут по-разному реализовывать однотипные конструкции.

\subsubsection{Повышение уровня абстракции}

Модель позволяет рассуждать о системе на более высоком уровне, чем исходный код.

Например, вместо описания SQL-команд, классов, контроллеров и сериализации можно описать сущности предметной области и связи между ними:

\[
    Customer \longleftrightarrow Order \longleftrightarrow Product.
\]

Затем генератор создаёт технические детали реализации.

\subsubsection{Упрощение сопровождения}

Если модель является главным источником истины, то изменение структуры системы можно выполнить в модели, а затем заново сгенерировать код.

Например, при добавлении поля:

\[
    User.phone: String
\]

генератор может автоматически обновить:

\begin{itemize}
    \item класс \texttt{User};
    \item DTO;
    \item форму ввода;
    \item миграцию БД;
    \item REST API;
    \item тестовые данные.
\end{itemize}

\subsubsection{Ускорение прототипирования}

Генерация кода полезна для быстрого создания прототипов.

Например, по модели предметной области можно быстро получить:

\begin{itemize}
    \item работающий CRUD-интерфейс;
    \item базовую БД;
    \item REST API;
    \item документацию API.
\end{itemize}

Это позволяет быстрее проверять гипотезы и согласовывать требования с заказчиком.

\subsubsection{Повторное использование знаний}

Правила генерации фиксируют архитектурные и технологические решения.

Например, если команда разработала удачный способ генерации REST-контроллеров, этот способ можно применять во многих проектах.

Таким образом, повторно используются не только библиотеки, но и архитектурные шаблоны.

\subsubsection{Повышение согласованности артефактов}

Одна модель может служить источником для нескольких артефактов:

\[
    M \longrightarrow \{Code, SQL, Tests, Documentation, API\ Spec\}.
\]

Например, из одной модели можно получить:

\begin{itemize}
    \item Java-классы;
    \item SQL-схему;
    \item OpenAPI-описание;
    \item HTML-документацию;
    \item тестовые данные.
\end{itemize}

Это снижает риск расхождения между документацией, кодом и базой данных.

\subsubsection{Повышение переносимости}

Если используется платформенно-независимая модель, то теоретически одну и ту же модель можно преобразовать в реализации для разных платформ:

\[
    PIM \longrightarrow PSM_{Java},
\]
\[
    PIM \longrightarrow PSM_{C\#},
\]
\[
    PIM \longrightarrow PSM_{Python}.
\]

Например, модель сущности может быть преобразована как в Java-класс, так и в C\#-класс.

\subsubsection{Поддержка формальной верификации}

Если модель имеет строгую семантику, то её можно анализировать до генерации кода.

Например, для конечного автомата можно проверить:

\begin{itemize}
    \item достижимость состояний;
    \item отсутствие тупиков;
    \item полноту переходов;
    \item отсутствие недетерминизма.
\end{itemize}

Это позволяет обнаруживать ошибки раньше, чем на этапе реализации.

\subsection{Недостатки автоматической генерации кода}

\subsubsection{Высокая стоимость начального внедрения}

Для эффективной генерации кода необходимо:

\begin{itemize}
    \item выбрать или разработать язык моделирования;
    \item определить метамодель;
    \item создать генератор;
    \item разработать шаблоны;
    \item настроить процесс сборки;
    \item обучить команду.
\end{itemize}

Эти затраты могут быть неоправданными для малых проектов.

\subsubsection{Сложность создания качественного генератора}

Генератор должен учитывать множество деталей:

\begin{itemize}
    \item особенности языка программирования;
    \item архитектурный стиль;
    \item обработку ошибок;
    \item безопасность;
    \item производительность;
    \item расширяемость;
    \item интеграцию с существующим кодом.
\end{itemize}

Плохо спроектированный генератор может создавать трудно сопровождаемый код.

\subsubsection{Ограниченность выразительности модели}

Модель обычно описывает только некоторый аспект системы.

Например, диаграмма классов хорошо описывает структуру данных, но плохо отражает:

\begin{itemize}
    \item сложные алгоритмы;
    \item оптимизации;
    \item конкурентность;
    \item нестандартную бизнес-логику;
    \item особенности интеграции с внешними системами.
\end{itemize}

Поэтому полностью сгенерировать сложную систему часто невозможно.

\subsubsection{Риск расхождения модели и кода}

Если разработчики изменяют сгенерированный код вручную, то может возникнуть рассогласование:

\[
    Model \neq Code.
\]

Например, в модели поле называется \texttt{email}, а в коде разработчик вручную переименовал его в \texttt{mail}. При следующей генерации изменения могут быть потеряны.

Для предотвращения этого необходимо строго определить источник истины:

\begin{itemize}
    \item либо модель является главным артефактом;
    \item либо код является главным артефактом;
    \item либо используется двусторонняя синхронизация.
\end{itemize}

\subsubsection{Проблема повторной генерации}

Если код был сгенерирован один раз, а затем изменён вручную, повторная генерация может перезаписать ручные изменения.

Пример опасной ситуации:

\begin{verbatim}
// generated code
public class UserService {
    public void save(User user) {
        // manually added validation
        validate(user);

        repository.save(user);
    }
}
\end{verbatim}

При повторной генерации метод может быть перезаписан:

\begin{verbatim}
public class UserService {
    public void save(User user) {
        repository.save(user);
    }
}
\end{verbatim}

В результате ручная валидация будет потеряна.

\subsubsection{Сложность отладки}

Если ошибка возникает в сгенерированном коде, то не всегда понятно, где её исправлять:

\begin{itemize}
    \item в модели;
    \item в шаблоне генератора;
    \item в самом генераторе;
    \item в ручном коде;
    \item в настройках трансформации.
\end{itemize}

Кроме того, сгенерированный код иногда бывает менее читаемым, чем написанный вручную.

\subsubsection{Зависимость от инструмента}

Использование генератора может привести к vendor lock-in, то есть зависимости от конкретного инструмента, платформы или формата модели.

Проблемы возникают, если:

\begin{itemize}
    \item инструмент перестаёт поддерживаться;
    \item формат модели закрыт;
    \item генератор невозможно адаптировать;
    \item отсутствует экспорт в стандартные форматы;
    \item трудно перенести проект на другую технологию.
\end{itemize}

\subsubsection{Снижение гибкости реализации}

Генератор обычно хорошо работает для типовых случаев, но плохо справляется с исключениями.

Например, если все сущности должны генерироваться как обычные CRUD-контроллеры, то нестандартная бизнес-логика может потребовать обходных решений.

Это приводит к появлению:

\begin{itemize}
    \item специальных флагов в модели;
    \item условных шаблонов;
    \item ручных вставок;
    \item нестандартных расширений;
    \item усложнения генератора.
\end{itemize}

\subsubsection{Сложность коллективной работы с моделями}

Если модель хранится в бинарном формате или в сложном XML-файле, то возникают трудности:

\begin{itemize}
    \item слияния изменений;
    \item просмотра различий;
    \item code review;
    \item разрешения конфликтов;
    \item версионирования.
\end{itemize}

Текстовые DSL-модели обычно лучше подходят для командной разработки.

\subsubsection{Риск генерации избыточного кода}

Генератор может создавать слишком много шаблонного кода:

\begin{itemize}
    \item неиспользуемые методы;
    \item лишние классы;
    \item чрезмерно сложную иерархию;
    \item избыточные конфигурации.
\end{itemize}

Это увеличивает размер проекта и усложняет сопровождение.

\subsubsection{Проблемы производительности}

Сгенерированный код не всегда оптимален.

Генератор может отдавать приоритет универсальности, а не эффективности. Например, он может создавать:

\begin{itemize}
    \item лишние уровни абстракции;
    \item неэффективные запросы к БД;
    \item избыточные преобразования данных;
    \item чрезмерное использование reflection;
    \item неоптимальные структуры данных.
\end{itemize}

Для критичных по производительности систем может потребоваться ручная оптимизация.

\subsection{Критерии применимости автоматической генерации кода}

Автоматическая генерация кода наиболее полезна, если:

\begin{enumerate}
    \item в проекте много повторяющегося кода;
    \item предметная область хорошо формализуется;
    \item архитектура достаточно стабильна;
    \item имеется набор типовых шаблонов реализации;
    \item команда готова поддерживать модели и генератор;
    \item важна согласованность документации, кода и схем данных;
    \item система относится к семейству похожих систем.
\end{enumerate}

Генерация кода менее целесообразна, если:

\begin{enumerate}
    \item проект небольшой;
    \item бизнес-логика сильно нестандартна;
    \item требования часто и хаотично меняются;
    \item нет ресурсов на поддержку генератора;
    \item команда не имеет опыта модельно-ориентированной разработки;
    \item сгенерированный код сложно интегрировать с существующей системой.
\end{enumerate}

\subsection{Сравнение ручного программирования и генерации кода}

\[
\begin{array}{|p{4cm}|p{5cm}|p{5cm}|}
\hline
\textbf{Критерий} & \textbf{Ручное программирование} & \textbf{Генерация по модели} \\
\hline
Скорость создания типового кода & Ниже & Выше \\
\hline
Гибкость & Высокая & Ограничена возможностями модели и генератора \\
\hline
Единообразие архитектуры & Зависит от дисциплины команды & Обеспечивается шаблонами \\
\hline
Начальные затраты & Ниже & Выше \\
\hline
Сопровождение типовых изменений & Может быть трудоёмким & Упрощается при корректной модели \\
\hline
Отладка & Обычно проще & Может быть сложнее из-за промежуточного уровня модели \\
\hline
Зависимость от инструментов & Ниже & Выше \\
\hline
Риск рассогласования артефактов & Выше & Ниже, если модель является источником истины \\
\hline
\end{array}
\]

\subsection{Пример: генерация SQL-схемы по модели}

Пусть задана модель сущности:

\[
\begin{array}{|c|}
\hline
\text{Book} \\
\hline
id: Integer \\
title: String \\
author: String \\
year: Integer \\
\hline
\end{array}
\]

Правила преобразования:

\[
    Integer \mapsto INTEGER,
\]
\[
    String \mapsto VARCHAR(255).
\]

Тогда результат генерации:

\begin{verbatim}
CREATE TABLE book (
    id INTEGER PRIMARY KEY,
    title VARCHAR(255),
    author VARCHAR(255),
    year INTEGER
);
\end{verbatim}

Алгоритм генерации SQL:

\begin{enumerate}
    \item Получить имя сущности.
    \item Преобразовать имя сущности в имя таблицы.
    \item Для каждого атрибута определить SQL-тип.
    \item Сформировать список столбцов.
    \item Добавить ограничения первичного ключа.
    \item Сформировать команду \texttt{CREATE TABLE}.
\end{enumerate}

\subsection{Пример: генерация кода конечного автомата}

Рассмотрим конечный автомат турникета:

\[
    S = \{Locked, Unlocked\},
\]
\[
    E = \{coin, push\}.
\]

Переходы:

\[
    Locked \xrightarrow{coin} Unlocked,
\]
\[
    Unlocked \xrightarrow{push} Locked.
\]

Таблица переходов:

\[
\begin{array}{|c|c|c|}
\hline
\textbf{Текущее состояние} & \textbf{Событие} & \textbf{Новое состояние} \\
\hline
Locked & coin & Unlocked \\
\hline
Unlocked & push & Locked \\
\hline
\end{array}
\]

Сгенерированный код на псевдо-Java:

\begin{verbatim}
enum State {
    LOCKED,
    UNLOCKED
}

class Turnstile {
    private State state = State.LOCKED;

    public void handle(Event event) {
        switch (state) {
            case LOCKED:
                if (event == Event.COIN) {
                    state = State.UNLOCKED;
                }
                break;

            case UNLOCKED:
                if (event == Event.PUSH) {
                    state = State.LOCKED;
                }
                break;
        }
    }
}
\end{verbatim}

Здесь модель автомата является более компактной и наглядной, чем код. Генератор автоматически создаёт перечисление состояний, обработчик событий и таблицу переходов.

\subsection{Типовые архитектурные приёмы при генерации кода}

\subsubsection{Generation Gap}

Шаблон ``generation gap'' разделяет сгенерированный и пользовательский код.

Например, генератор создаёт класс:

\begin{verbatim}
public abstract class UserBase {
    protected String name;
    protected String email;
}
\end{verbatim}

Разработчик пишет наследника:

\begin{verbatim}
public class User extends UserBase {
    public boolean hasValidEmail() {
        return email != null && email.contains("@");
    }
}
\end{verbatim}

При повторной генерации изменяется только \texttt{UserBase}, а пользовательский класс \texttt{User} сохраняется.

\subsubsection{Protected Regions}

В сгенерированном коде выделяются области, которые генератор не перезаписывает:

\begin{verbatim}
public class UserService {
    // <protected>
    public void customValidation(User user) {
        // ручной код
    }
    // </protected>
}
\end{verbatim}

Недостаток этого подхода состоит в том, что сгенерированный и ручной код смешиваются в одном файле.

\subsubsection{Partial Classes}

Если язык поддерживает частичные классы, сгенерированная и ручная части могут находиться в разных файлах.

Например, в C\#:

\begin{verbatim}
// User.generated.cs
public partial class User {
    public string Name { get; set; }
}

// User.cs
public partial class User {
    public bool IsValid() {
        return !string.IsNullOrEmpty(Name);
    }
}
\end{verbatim}

\subsection{Качество сгенерированного кода}

Качество сгенерированного кода оценивается по следующим критериям:

\begin{enumerate}
    \item \textbf{Корректность} --- код должен соответствовать модели.
    \item \textbf{Читаемость} --- код должен быть понятен разработчику.
    \item \textbf{Поддерживаемость} --- код должен допускать изменения.
    \item \textbf{Идиоматичность} --- код должен соответствовать стилю целевого языка.
    \item \textbf{Производительность} --- код не должен содержать существенных неоправданных накладных расходов.
    \item \textbf{Тестируемость} --- код должен быть удобен для модульного и интеграционного тестирования.
    \item \textbf{Безопасность} --- генератор не должен создавать уязвимые конструкции.
\end{enumerate}

\subsection{Типичные ошибки при использовании генерации кода}

\begin{enumerate}
    \item Генерация кода без чётко определённой метамодели.
    \item Использование модели только как картинки, не имеющей строгой семантики.
    \item Ручное изменение сгенерированного кода без стратегии сохранения изменений.
    \item Чрезмерная генерация всего подряд.
    \item Отсутствие тестов для генератора.
    \item Сильная зависимость от конкретного инструмента.
    \item Отсутствие версии модели в системе контроля версий.
    \item Сложные шаблоны, которые трудно сопровождать.
    \item Попытка заменить моделями всё программирование.
\end{enumerate}

\subsection{Практические рекомендации}

\begin{enumerate}
    \item Генерировать прежде всего повторяющийся и формализуемый код.
    \item Не смешивать ручной и сгенерированный код без необходимости.
    \item Хранить модель в системе контроля версий.
    \item Использовать текстовые DSL, если важны code review и удобное сравнение изменений.
    \item Делать генератор детерминированным: одна и та же модель должна давать один и тот же результат.
    \item Покрывать генератор тестами.
    \item Документировать правила генерации.
    \item Определить, что является главным источником истины: модель или код.
    \item Предусмотреть механизм расширения сгенерированного кода.
    \item Не использовать генерацию там, где ручной код проще и понятнее.
\end{enumerate}

\subsection{Итоговая оценка}

Автоматическая генерация кода по модели является мощным средством повышения производительности и согласованности разработки. Она особенно эффективна в задачах, где присутствует большое количество однотипного кода и хорошо формализуемая предметная область.

Главные достоинства подхода:

\begin{itemize}
    \item ускорение разработки;
    \item снижение числа рутинных ошибок;
    \item единообразие архитектуры;
    \item повышение уровня абстракции;
    \item согласованность кода, документации и схем данных;
    \item возможность повторного использования архитектурных решений.
\end{itemize}

Главные недостатки:

\begin{itemize}
    \item высокая стоимость внедрения;
    \item сложность генераторов;
    \item ограниченность моделей;
    \item риск рассогласования модели и кода;
    \item сложность отладки;
    \item зависимость от инструментов;
    \item возможная избыточность и неоптимальность сгенерированного кода.
\end{itemize}

Таким образом, генерация кода по модели целесообразна не как универсальная замена программированию, а как инженерный инструмент автоматизации повторяющихся и формализуемых частей разработки.

\end{document}
